\scanpage{81}
\mspage{36}

\deleted{Le cas ou $V$}

Posons, pour une bande $J$
\[
\begin{aligned}
(62)\qquad V_J&=\text{provoisinage tubulaire de }X_J^*\text{ dans }X\\
&=\bigcup_{i\in J}V_i=\varinjlim_{i,j\in J}(V_i,V_{ij})
\end{aligned}
\]
On s'intéresse aux protopos de la forme
\[
(63)\qquad V_J^L\overset{\mathrm{def}}{=}V_J\cap X_L^*,
\]
\margin{\uline{Ex} $J\subset\deleted{J'}\inserted{K\ \text{(idéal de $J'$)}}$ inclusion de bandes, \deleted{\unsure{disposition}} $X_J^*\subset X_{J'}^*$, on prend
$V(X_J^*,X_{J'}^*)\setminus V(X_J^*,X_{J'}^*)\cap X_K^*$}
\margin{protopos de raccord entre $X_J^*$ et $X_L^*$
\diagram{\begin{tikzcd}[column sep=small,row sep=small]
V_J^L\arrow[r,hook]\arrow[d,hook]&X_L^*\arrow[r,hook]&V_L\\
V_J&&
\end{tikzcd}}
\deleted{\unsure{NB} $V_J^{\unsure{L}}\cap V_{J'}^{L'}=V_{J\cap J'}^{L\cap L'}$}}
où $J,L$ sont deux bandes. Par exemple on retrouve les $X_J^*$ comme
\[
(64)\qquad V_J^J=X_J^*
\]
Soit $L_0\subset L$ \unsure{l'ens|) l'ens} des $l\in L$ qui majorent un \unsure{élém|élément} de $J$, et $J_0$ l'ens des $j\in J$ qui mino-\\rent un élément de $L$, on a donc
\[
(65)\qquad V_J^L=V_J^{L_0}=V_{J_0}^L=V_{J_0}^{L_0}
\]
où $J_0,L_0$ sont des sous-bandes de $J,L$ respectivement ($J_0$ est = un idéal de $J$, $L_0$ un coidéal de $L$) telles que tt élément de $J_0$ est majoré par un élément de $L_0$, tt élément de $L_0$ majore un élément de $J_0$. On est ainsi ramené au cas de bandes dans la relation de domination
\[
(66)\qquad J\prec L
\]
\margin{qui signifie que tt élément de $J$ majoré par un élément de $L$, tt \unsure{élément|élémt} de $L$ majore un él. de $J$.}
\unsure{de|donc}
% The light, thin bottom-left foliation is in another hand.
\margin{80}
